\section{Experiments}

\paragraph{Comparison to state-of-the-art real-time object detectors.} Table~\ref{tab:XLMS} summarizes the performance of DEIMv2 across the S, M, L, and X variants, demonstrating substantial improvements over prior state-of-the-art detectors. For instance, the largest variant, DEIMv2-X, attains 57.8 AP with only $\sim$50M parameters and 151 GFLOPs, surpassing the previous best DEIM-X (56.5 AP with 62M parameters and 202 GFLOPs). This demonstrates that DEIMv2 can deliver superior accuracy with both fewer parameters and lower computational cost. At the smaller end, DEIMv2-S sets a new milestone as the first model with fewer than 10M parameters to exceed the 50 AP threshold on COCO, achieving 50.9 AP with only 11M parameters and 26 GFLOPs. This marks a clear improvement over the prior DEIM-S (49.0 AP with 10M parameters), while requiring nearly the same model size. While CNN-based backbones are generally more hardware-friendly, our ViT-based backbone achieves a lightweight design with fewer parameters and lower FLOPs, offering better scalability and deployment flexibility. It is worth noting that the latency of the proposed methods has not been optimized. Techniques such as Flash Attention~\cite{dao2022flashattention}, as in Yolov12~\cite{tian2025yolov12}, could further accelerate inference. Overall, the reduced FLOPs highlight the potential of ViT-based backbones to achieve low-latency performance with proper optimization.%

Interestingly, when comparing DINOv3-based DEIMv2 models with their previous DEIM counterparts under comparable parameter and FLOP budgets, the accuracy gains primarily arise from improvements on medium and large objects, while performance on small objects remains largely unchanged. For instance, DEIMv2-S achieves 55.3 AP$_M$ and 70.3 AP$_L$, clearly surpassing DEIM-S (52.6 AP$_M$ and 65.7 AP$_L$), yet the small-object scores are nearly identical (31.4 vs. 30.4 AP$_S$). A similar trend is observed for larger models: DEIMv2-X improves AP$_M$ from 61.4 to 62.8 and AP$_L$ from 74.2 to 75.9, while its small-object AP (39.2) remains close to that of DEIM-M (38.8). These results indicate that DEIMv2’s main advantage lies in enhancing the representation and detection of medium-to-large objects, whereas small-object detection remains a challenge across scales. This observation further confirms that DINOv3 excels at capturing strong global semantics but has limited ability to represent fine-grained details. Exploring ways to better integrate DINOv3 features into real-time detectors thus represents an interesting direction for future work. 


\paragraph{Comparison to competitive ultra-light object detectors.} The ultra-light variants of DEIMv2 also exhibit strong performance, as summarized in Table~\ref{tab:my_sorted_label_asc}. DEIMv2-Atto, with only 0.49M parameters, achieves performance comparable to NanoDet-M despite its substantially smaller size. Similarly, DEIMv2-Pico attains performance on par with YOLOv10-N~\cite{wang2024yolov10} while requiring less than half the parameters. These results underscore the effectiveness of DEIMv2 in extremely compact regimes and highlight its suitability for deployment on resource-constrained edge devices.



